\documentclass[tikz,border=8pt]{standalone}
\usepackage{ctex}
\usepackage{amsmath}
\usetikzlibrary{calc, arrows.meta, positioning}
\begin{document}
\begin{tikzpicture}[thick,
  ax/.style={->, semithick, black},
  curve/.style={blue!70!black, ultra thick},
]
  % ============ 总标题 ============
  \node[font=\bfseries] at (4.5,7.0)
    {CDF 统一刻画分布：离散情形 PMF $=$ 跳跃高度，连续情形 PDF $=$ 斜率};

  % =================================================================
  % 面板 1（左上）：离散型 CDF（阶梯）  —— P(0)=1/4,P(1)=1/2,P(2)=1/4
  % =================================================================
  \begin{scope}[shift={(0,4)}]
    \node[font=\small\bfseries] at (1.5,2.65) {离散型 CDF $F(x)$（阶梯函数）};
    \draw[ax] (-0.3,0) -- (3.4,0) node[right,font=\scriptsize] {$x$};
    \draw[ax] (0,-0.3) -- (0,2.2) node[above,font=\scriptsize] {$F$};
    \node[font=\scriptsize,left] at (-0.05,1.8) {1};
    \foreach \x in {1,2} {\draw (\x,-0.06)--(\x,0.06); \node[font=\scriptsize,below] at (\x,-0.05){\x};}
    \node[font=\scriptsize,below] at (0,-0.05) {0};
    % 阶梯段（prob ×1.8）
    \draw[curve] (-0.3,0) -- (0,0);
    \draw[curve] (0,0.45) -- (1,0.45);
    \draw[curve] (1,1.35) -- (2,1.35);
    \draw[curve] (2,1.8) -- (3.4,1.8);
    % 实心/空心端点
    \foreach \x/\y in {0/0.45,1/1.35,2/1.8} {\fill[blue!70!black] (\x,\y) circle (2pt);}
    \foreach \x/\y in {0/0,1/0.45,2/1.35} {\draw[blue!70!black,fill=white] (\x,\y) circle (2pt);}
    % 标注：x=1 处跳跃高度 = p(1)
    \draw[<->,red,thin] (1.08,0.45) -- (1.08,1.35) node[midway,right,font=\scriptsize,red] {跳跃 $=p(1)=\tfrac12$};
  \end{scope}

  % =================================================================
  % 面板 2（右上）：离散型 PMF（跳跃高度）
  % =================================================================
  \begin{scope}[shift={(6,4)}]
    \node[font=\small\bfseries] at (1.5,2.65) {离散型 PMF $p(x)$（即各跳跃高度）};
    \draw[ax] (-0.3,0) -- (3.4,0) node[right,font=\scriptsize] {$x$};
    \draw[ax] (0,-0.3) -- (0,2.2) node[above,font=\scriptsize] {$p$};
    \foreach \x in {0,1,2} {\draw (\x,-0.06)--(\x,0.06); \node[font=\scriptsize,below] at (\x,-0.05){\x};}
    % 竖直「砝码」stems
    \draw[red,ultra thick,-{Stealth[length=2mm]}] (0,0) -- (0,0.45) node[above,font=\scriptsize,black]{$\tfrac14$};
    \draw[red,ultra thick,-{Stealth[length=2mm]}] (1,0) -- (1,0.9)  node[above,font=\scriptsize,black]{$\tfrac12$};
    \draw[red,ultra thick,-{Stealth[length=2mm]}] (2,0) -- (2,0.45) node[above,font=\scriptsize,black]{$\tfrac14$};
  \end{scope}

  % =================================================================
  % 面板 3（左下）：连续型 CDF（光滑 S 曲线）
  % =================================================================
  \begin{scope}[shift={(0,0)}]
    \node[font=\small\bfseries] at (1.5,2.65) {连续型 CDF $F(x)$（光滑递增）};
    \draw[ax] (-0.3,0) -- (3.4,0) node[right,font=\scriptsize] {$x$};
    \draw[ax] (0,-0.3) -- (0,2.2) node[above,font=\scriptsize] {$F$};
    \node[font=\scriptsize,left] at (-0.05,1.8) {1};
    \draw[curve,smooth,domain=0:3,samples=60] plot (\x,{1.8/(1+exp(-3*(\x-1.5)))});
    % 拐点处切线，斜率 = f(x)
    \draw[red,thin] (1.0,0.225) -- (2.0,1.575);
    \node[font=\scriptsize,red,anchor=west] at (2.05,1.5) {斜率 $=f(x)$};
    \fill[red] (1.5,0.9) circle (1.6pt);
  \end{scope}

  % =================================================================
  % 面板 4（右下）：连续型 PDF（钟形密度）
  % =================================================================
  \begin{scope}[shift={(6,0)}]
    \node[font=\small\bfseries] at (1.5,2.65) {连续型 PDF $f(x)$（即 CDF 的斜率）};
    \draw[ax] (-0.3,0) -- (3.4,0) node[right,font=\scriptsize] {$x$};
    \draw[ax] (0,-0.3) -- (0,2.2) node[above,font=\scriptsize] {$f$};
    \draw[curve,fill=blue!12,smooth,domain=0:3,samples=70]
      plot (\x,{1.6*exp(-(\x-1.5)^2/0.28)}) -- (3,0) -- (0,0) -- cycle;
    % 半透明注释，避开峰值（峰在 x=1.5），放右侧
    \node[font=\scriptsize, fill=yellow!18, draw=yellow!60!black, rounded corners=2pt,
          align=center, fill opacity=0.85, text opacity=1, anchor=east]
      at (3.35,1.55) {面积 $=$\\概率};
  \end{scope}

\end{tikzpicture}
\end{document}
